\section{Discussion}
\label{sec:discussion}

\paragraph{The answer to the motivating question is no---with a precise
scope.} On this large, noisy financial dataset, under a pre-registered
protocol with matched compute budgets, the Spectral Optimizer did not
reduce the effect of noise: there is no evidence of benefit under the
affordable tuning budget, and at the pre-registered operating point the
filter significantly degraded out-of-sample performance on both an MLP and
a GRU, erasing $\sim$82--83\% of a small ($+0.0064$) baseline edge. The
result survived an adversarial audit that re-executed the experiment on
fresh seeds and an independently constructed evaluation split. The claim is
scoped to a subsampled, low-capacity pipeline whose verdict-block edge sits
below the practitioner sanity band, and to a single spectral operating
point.

\paragraph{Why the label-noise story does not transfer.} The two
mechanistic results give a consistent account. By
Theorem~\ref{thm:target-independence}, on scalar-output MSE the quantity
the filter thresholds---the eigenspectrum of the normalized per-sample
gradient similarity matrix---is exactly target-independent at fixed
parameters. ``Directions many samples agree on'' are, on this task class,
directions of the \emph{feature/Jacobian factor structure}, not of target
signal; the coherence-amplifier interpretation from the label-noise setting
has nothing to grip. Cross-entropy classification escapes the theorem
(per-sample gradients are not scalar multiples of a target-independent
vector), which explains why the same diagnostics were informative on MNIST
label noise and are not here. Empirically, the C4 control closes the loop
from the other end: conditioning on the kept-$k$ and update-norm
trajectories, replacing MP-eigenselection with a random sample-subspace
changes nothing detectable ($\pm 0.004$ resolution), on either
architecture, in either engagement regime ($k \approx 1$ with
10$\times$ amplification on the MLP; $k \approx 60$ with mild attenuation
on the GRU).

\paragraph{The harm is a class effect, not a spectral effect.} All three
agreement-style arms---spectral, random-subspace, and sign-agreement---hurt
by similar amounts. Whatever damage is done appears attributable to
per-sample-agreement-style update restriction/rescaling itself. Two
candidate stories for \emph{why} such interventions hurt here are already
constrained by the data: it is not a stability trade (corr-Sharpe fell,
robustly so on the GRU), and it is not a Feldman-style tail effect (the
unconditioned dispersion split shows uniform harm, not harm concentrated in
high-dispersion eras). A plausible remaining account---that in a regime
where the mean gradient is the only weakly informative direction, replacing
it with an amplified low-rank projection of factor structure simply
injects variance where there is no exploitable coherent signal---is
consistent with the prior project's own weak-signal failure mode, but our
design does not isolate it; we flag it as interpretation, not finding.

\paragraph{On the strength of the negative.} The pre-registered machinery
was designed so that a negative would be meaningful: the F4 power gate
showed all three verdict categories reachable (power 0.91 at the frozen
threshold) before any slot was spent; engagement diagnostics confirm the
mechanism did sustained, selective work during every verdict run; and the
audit classified the outcome as an adequately powered honest negative
(exit reason: \emph{true-null}) rather than an artifact. The main
interpretive guard is deliberate and one-sided: because the spectral arm
afforded only $\sim$1 compute-matched tuning trial, ``hurts'' cannot be
distinguished from ``under-tuned spectral arm'' at this budget, so the
paper's claim stays at ``no evidence of benefit under the affordable
tuning budget'' with the raw numbers alongside. The mirrored signature
(the baseline's regularization grid-boundary trend) cuts the other way:
a better-tuned baseline would, if anything, widen the gap.

\paragraph{What would change our mind.} A demonstration that a tuned
spectral configuration (wider operating-point sweep at a declared,
enlarged, matched budget) recovers or beats the baseline in this regime;
or a full-scale (non-subsampled) evaluation in which the verdict-block
edge returns to the practitioner band and the effect reverses. Both are
concrete, costed follow-ups in Section~\ref{sec:future-work}.
